\documentclass[11pt,a4paper]{amsart}
% Adapted from an OpenAI Math manuscript. See PROVENANCE.md.

\usepackage{cgeometry}
\crefname{maintheorem}{Theorem}{Theorems}
\Crefname{maintheorem}{Theorem}{Theorems}

\addbibresource{ComplexGeometry.bib}
\defbibheading{bibliography}[\refname]{\section*{#1}}
\setlength{\bibitemsep}{0pt}
\renewcommand*{\bibfont}{\footnotesize}

\allowdisplaybreaks[2]
\emergencystretch=1em

%%%%%%%%%%%%%%%%%%%%%%%%%%%%%%%%%%%%%%%%%%%%%%%%%%%%%%%%%%%%%%%%%%%%%%%%%%%%%
%% Notation
%%%%%%%%%%%%%%%%%%%%%%%%%%%%%%%%%%%%%%%%%%%%%%%%%%%%%%%%%%%%%%%%%%%%%%%%%%%%%


\usepackage{xurl}
\usetikzlibrary{arrows.meta,calc,positioning,patterns}

\DeclareMathOperator{\mult}{mult}
\DeclareMathOperator{\conv}{conv}
\DeclareMathOperator{\length}{length}

\title{Maximal multipoint Seshadri constants in positive characteristic}
\author{}
\date{}

\begin{document}

\begin{abstract}
Let $L$ be an ample line bundle on a smooth integral projective variety $X$ of dimension $n\geq2$ over an algebraically closed field of positive characteristic. We prove that its multipoint Seshadri constant at the geometric generic tuple of $r$ points is $(L^n/r)^{1/n}$ for all sufficiently large $r$. The construction uses Hasse derivatives and flat specialization of finite jet schemes.
\end{abstract}

\maketitle

% Adapted from an OpenAI Math manuscript. See PROVENANCE.md.
\section{Introduction}\label{sec:introduction}

Let $k$ be an algebraically closed field of characteristic $p>0$, let $X/k$ be a smooth integral projective variety of dimension $n$, and let $L$ be an ample line bundle. For $r\ge1$, write
\[
 U_r=\{(x_1,\ldots,x_r)\in X^r:x_i\ne x_j\text{ for }i\ne j\}
\]
for the open configuration space of ordered distinct points. Put $K_r=\overline{k(U_r)}$, and let $p_1,\ldots,p_r\in X(K_r)$ be the points specified by its geometric generic point. Define
\begin{equation}\label{eq:generic-seshadri}
 \varepsilon_r(X,L)=\inf_C
 \frac{L_{K_r}\cdot C}{\sum_{i=1}^r\mult_{p_i}C},
\end{equation}
where $C$ ranges over integral curves in $X_{K_r}$ meeting at least one $p_i$. Multiplicity at a point outside $C$ is zero. Throughout, we use this ordinary, curve-defined Seshadri constant.

\begin{maintheorem}\label{thm:main}
Let $X/k$ be a smooth integral projective variety over an algebraically closed field of characteristic $p>0$, with $\dim X=n\ge2$, and let $L$ be ample. There is an integer $r_0=r_0(X,L)$ such that
\[
 \varepsilon_r(X,L)=\left(\frac{L^n}{r}\right)^{1/n}
 \qquad\text{for every integer }r\ge r_0.
\]
\end{maintheorem}

This exposition follows the proof in \cite{OpenAIMathSource}.


The geometric generic formulation works over countable algebraically closed fields. On the blowup of the geometric generic tuple, the theorem gives the nef class
\[
 \pi^*L_{K_r}-(L^n/r)^{1/n}\sum_{i=1}^r E_i
\]
of top self-intersection zero, where the $E_i$ are the exceptional divisors.

The proof adapts the complete-intersection flag and simplex reshaping of the complex case \cite[Sections~2--6]{OpenAI2026Higher}. In positive characteristic, formal substitution and reduction of jet matrices over $k[[s]]$ replace analytic rescaling. The equation $y^2+xy-x^3=0$ has two distinct tangent lines at the node, including in characteristic two. The branch estimate provides the two-coordinate reshaping, with jet orders expressed through powers of maximal ideals (\cref{sec:transfer,sec:flag,sec:reshape}).

For fixed sufficiently large $r$, put $w=(L^n/r)^{1/n}$. The reshaped monomial spaces have supports in dilates of a simplex with intercepts $w,\ldots,w,rw$. Their graded addition law fills compact interior regions in large degree. The Chinese remainder theorem along the last coordinate gives jet separation up to order $\lfloor N\theta\rfloor$ for every fixed $0<\theta<w$ in large degree $N$ (\cref{sec:interpolation}). For each degree, the transferred rank condition holds at the geometric generic tuple because it holds on a nonempty open subset. \Cref{sec:positivity} turns these jet orders into intersection bounds and lets $\theta$ increase to $w$.

% Adapted from an OpenAI Math manuscript. See PROVENANCE.md.
\section{Leading terms and transfer of jets}\label{sec:transfer}

We first turn spaces of sections into finite sets of exponents.  The essential point is that jet separation by the resulting monomials implies jet separation by the original sections.  We will also make two elementary changes to the exponent sets while retaining this implication. These operations follow the support method of \cite[Section~2]{OpenAI2026Higher}; the formal transfer below makes the method available in arbitrary characteristic.

Throughout, $k$ is algebraically closed.  The arguments of this section work in every characteristic.  For a smooth $d$-dimensional variety $Y$, a line bundle $A$, and $y\in Y(k)$, the jets through order $m\geq0$ are
\[
 J_y^m(A)=A_y/\mathfrak{m}_y^{m+1}A_y.
\]
A finite-dimensional space $W\subset H^0(Y,A)$ \emph{separates these jets at} $y_1,\ldots,y_r$ if the evaluation map $W\to\bigoplus_i J_{y_i}^m(A)$ is surjective.  The locus of such ordered distinct tuples is open in $Y^r$: the jet spaces are the fibers of the locally free sheaf of principal parts of $A$, of rank $\binom{d+m}{d}$, and surjectivity is a rank condition.  The description of principal parts by the $m$th infinitesimal neighborhood of the diagonal also shows compatibility with field extension.  Local coordinates for a smooth variety identify its fibers with truncated polynomial rings. For the diagonal description and base change, see \cite[Tags~\href{https://stacks.math.columbia.edu/tag/0H90}{0H90} and~\href{https://stacks.math.columbia.edu/tag/0H8Z}{0H8Z}]{Stacks}.  Local freeness follows as well from the filtration whose graded pieces are $A\otimes\operatorname{Sym}^j\Omega_Y^1$, $0\leq j\leq m$.

Put $\mathbb{T}^d=(\mathbb{G}_{m,k})^d$.  For a finite set $S\subset\mathbb{Z}^d$, write
\[
 M(S)=\operatorname{span}_k\{x^\alpha:\alpha\in S\}
 \subset k[x_1^{\pm1},\ldots,x_d^{\pm1}].
\]
Here and below, auxiliary existence statements about points on a torus refer to $k$-points.

\begin{definition}\label{def:graded}
A family of finite sets $(S_N)_{N\geq1}$ in $\mathbb{Z}^d$ is \emph{graded} if
\[
 S_N+S_{N'}\subset S_{N+N'}\qquad(N,N'\geq1).
\]
A transformation $(S_N)\mapsto(T_N)$ has \emph{backward jet transfer} if, for every $N,r\geq1$ and $m\geq0$, separation of jets through order $m$ by $M(T_N)$ at some $r$ distinct torus points implies the same assertion for $M(S_N)$, possibly at a different tuple.
\end{definition}

\subsection{Least exponents}

A monomial order on $\mathbb{Z}_{\geq0}^h$ means an addition-compatible total well-order.  For a nonzero formal series $f\in k[[z_1,\ldots,z_h]]$, let $\nu(f)$ be the least exponent occurring in $f$, for a fixed such order.  In particular,
\begin{equation}\label{eq:least-product}
 \nu(fg)=\nu(f)+\nu(g).
\end{equation}
Indeed, the product of the two least terms is the unique least term in the product.

\begin{lemma}\label{lem:least-terms}
For a finite-dimensional subspace $W\subset k[[z_1,\ldots,z_h]]$, the set $\{\nu(f):0\neq f\in W\}$ has cardinality $\dim_k W$. More generally, suppose finite-dimensional spaces $W_{N,\gamma}$, with only finitely many nonzero spaces in each degree, are indexed by $N\geq1$ and $\gamma\in\mathbb{Z}^{d-h}$ and satisfy
\[
 W_{N,\gamma}W_{N',\gamma'}
 \subset W_{N+N',\gamma+\gamma'}.
\]
Then recording the pairs $(\nu(f),\gamma)$ for $0\neq f\in W_{N,\gamma}$ gives a graded family, with cardinality $\sum_\gamma\dim_k W_{N,\gamma}$ in degree $N$.
\end{lemma}

\begin{proof}
Choose the smallest exponent occurring as $\nu(f)$ for a nonzero member of $W$, and choose a member whose coefficient there is $1$.  The kernel of that coefficient functional has codimension one.  Repeating in this kernel gives a basis with distinct least exponents.  Every nonzero linear combination has as its least exponent the smallest least exponent among the basis vectors with nonzero coefficients.  This proves the first assertion; the second follows label by label and from \eqref{eq:least-product}.
\end{proof}

For the transfer of jets we need to realize finitely many chosen least terms by a single positive weight vector.  The weights may depend on the degree; the monomial order defining the graded family will remain fixed.

\begin{lemma}\label{lem:exposing-weights}
Let $f_1,\ldots,f_q$ be nonzero formal series, and let $e_j=\nu(f_j)$. There is a vector $\lambda\in\mathbb{Z}_{>0}^h$ for which $e_j$ is the unique exponent of least $\lambda$-weight in $f_j$, simultaneously for all $j$, in either of the following cases:
\begin{enumerate}[label=\textup{(\roman*)}]
\item the order is lexicographic;
\item $h=2$ and, for a fixed real $t>0$, the order compares
$(\alpha+t\beta,\beta)$ lexicographically.
\end{enumerate}
\end{lemma}

\begin{proof}
In the lexicographic case set $\lambda_h=1$, and successively choose
\[
 \lambda_i>\max_j\sum_{a>i}\lambda_a(e_j)_a
 \qquad(i=h-1,\ldots,1).
\]
If an exponent is lexicographically greater than $e_j$, its gain at the first differing coordinate outweighs every possible loss in the later coordinates.

For (ii), the indicated order is a monomial order because every bounded weight range contains only finitely many nonnegative integral exponents. Choose a real bound larger than the weights of all the $e_j$.  Among the finitely many exponents below that bound, replacing $t$ by a sufficiently close rational $t'>t$ preserves all strict comparisons with the $e_j$ and resolves equal-weight comparisons by increasing $\beta$.  Exponents above the bound cannot interfere: their weights only increase, while the new weights of the $e_j$ stay below the bound.  Clearing the denominator of $(1,t')$ gives the required positive integral weights.
\end{proof}

\subsection{Formal specialization of jet matrices}

The next lemma supplies the passage from formal leading terms to actual points.  Ito's weighted degeneration argument gives such a jet-rank comparison over $\mathbb{C}$ \cite[Proposition~5.7]{Ito2013}.  Here substitution in truncated rings supplies a formal version in arbitrary characteristic, without differentiation or factorial denominators.

\begin{lemma}\label{lem:formal-transfer}
Let $1\leq h\leq d$, let $Y$ be a smooth integral $h$-dimensional $k$-variety, and choose $y\in Y(k)$ with formal coordinates $z_1,\ldots,z_h$.  Let $f_1,\ldots,f_q$ be regular functions on a common neighborhood of $y$, and let $\gamma_j\in\mathbb{Z}^{d-h}$. Suppose the expansion of $f_j$ has a term $c_jz^{e_j}$, with $c_j\neq0$, which is uniquely lowest for one common weight vector $\lambda\in\mathbb{Z}_{>0}^h$.

If the span of the monomials $Z^{e_j}u^{\gamma_j}$ separates jets through order $m$ at $r$ distinct points of $\mathbb{T}^h\times\mathbb{T}^{d-h}$, then the span of $f_ju^{\gamma_j}$ separates the same jets at $r$ distinct points in the product of that neighborhood with $\mathbb{T}^{d-h}$. The points may be chosen in any prescribed dense open subset of this product.  The statement also holds for sections written in a local frame of a line bundle.
\end{lemma}

\begin{proof}
Let $(Q_i,U_i)$, $1\leq i\leq r$, be a tuple at which the monomials separate the stated jets.  Shrink the neighborhood of $y$ to an affine one, denoted $Y^\circ$, on which all the functions are defined.  Put
\[
 R=k[[s]],\qquad F=k((s)),\qquad
 B=R[\delta_1,\ldots,\delta_d]/(\delta_1,\ldots,\delta_d)^{m+1},
\]
and write $\delta=(\delta',\delta'')$ according to the two factors. For each $i$, substitute
\begin{equation}\label{eq:formal-substitution}
 z_a=s^{\lambda_a}(Q_{i,a}+\delta'_a),\qquad
 u_b=U_{i,b}+\delta''_b.
\end{equation}
The first substitution is defined on the completed local ring, since $B$ is $s$-adically complete and every $\lambda_a$ is positive.  The second is defined on Laurent polynomials, since each $U_{i,b}$ is nonzero.  In particular, these substitutions define ring maps from the coordinate ring of $Y^\circ\times\mathbb{T}^{d-h}$ to $B$.

Use the classes of $\delta^\alpha$, $|\alpha|\leq m$, as an $R$-basis of $B$.  The images of $f_ju^{\gamma_j}$ under all $r$ substitutions give the $j$th column of a matrix with $r\binom{d+m}{d}$ rows.  Divide this column by $c_js^{\lambda\cdot e_j}$.  These normalized columns still belong to $B^{\oplus r}$, and their reductions modulo $s$ are precisely the jets of $Z^{e_j}u^{\gamma_j}$ at $(Q_i,U_i)$ in translated coordinates. The assumed surjectivity supplies a maximal minor nonzero modulo $s$. The original columns therefore span $B[1/s]^{\oplus r}$ over $F$.

We now interpret this rank statement as jet separation on a variety. Setting $\delta=0$ in \eqref{eq:formal-substitution} gives actual $F$-points $P_i$ of $Y^\circ\times\mathbb{T}^{d-h}$.  After inverting $s$, the corresponding ring map factors through the order-$m$ jet algebra of $(Y^\circ\times\mathbb{T}^{d-h})_F$ at $P_i$: elements nonvanishing at $P_i$ map to units, and its maximal ideal maps into $(\delta)$.  Both this jet algebra and $B[1/s]$ have dimension $\binom{d+m}{d}$ over $F$.  Surjectivity onto the direct sum of the latter spaces forces each induced map from a jet algebra to be an isomorphism.  It follows that the original functions separate the order-$m$ jets at the $P_i$ themselves.  These points are distinct, since two repeated centers would make two target blocks factor through one jet algebra, contradicting surjectivity.

Thus the open jet-separation locus has an $F$-point and is nonempty. The configuration space is irreducible, so this locus meets the configuration space of any prescribed dense open subset.  A nonempty open subset of a variety over the algebraically closed field $k$ has a $k$-point.  This proves the assertion.  A local frame reduces the line-bundle version to the one just proved.
\end{proof}

Together, Lemmas~\ref{lem:least-terms}--\ref{lem:formal-transfer} allow us to replace finite-dimensional spaces by their least exponents, using either of the two orders above, while preserving dimension and transferring jet separation back to the original spaces.  Labels record the exponents of variables that are left unchanged.

\subsection{Two elementary transformations}

\begin{lemma}\label{lem:unimodular}
For $A\in\operatorname{GL}_d(\mathbb{Z})$, the transformation $S_N\mapsto A S_N$ preserves cardinalities and gradedness, and has backward jet transfer.
\end{lemma}

\begin{proof}
An invertible integral linear map preserves cardinalities and addition. Its action on characters is induced by a torus automorphism, which identifies the corresponding jet evaluations.
\end{proof}

The other transformation uses orders of vanishing to place each row in an interval starting at zero.  The resulting exponents need not be consecutive.

\begin{lemma}\label{lem:compression}
Let $(S_N)$ be a graded family in $\mathbb{Z}^d$, and select one coordinate, written $z$.  For each degree $N$ and each label $\gamma\in\mathbb{Z}^{d-1}$, form its row space
\[
 W_{N,\gamma}
 =\operatorname{span}_k\{z^a:(a,\gamma)\in S_N\}.
\]
Replace this row by the set of orders at $z=1$ of nonzero elements of $W_{N,\gamma}$, retaining the label $\gamma$.  The resulting family preserves cardinalities and gradedness, and has backward jet transfer. If the original row lies in a real interval of length $b$, every new exponent lies in $[0,b]$.
\end{lemma}

\begin{proof}
Expand the row spaces in the parameter $z-1$.  \cref{lem:least-terms} proves the cardinality and gradedness assertions, because multiplication adds both degrees and labels.  \cref{lem:formal-transfer}, with $Y=\mathbb{G}_m$, $y=1$, and weight $1$, proves backward jet transfer.

For a nonempty row let $a_-$ and $a_+$ be its smallest and largest exponents.  Multiplication by $z^{-a_-}$ preserves orders at $1$ and turns each member into a polynomial of degree at most $a_+-a_-\leq b$. Its order of vanishing is at most its degree, proving the bound.
\end{proof}

Consequently, suppose a family is bounded by dilates of a polygon. Replacing each slice of that polygon parallel to the compressed coordinate by an interval of the same length starting at zero gives a new region whose dilates contain the compressed family.  This applies equally to real slice endpoints and to every dilation of the polygon.  In characteristic $p$, a compressed row need not contain all intermediate integers: for example, $\operatorname{span}_k\{1,z^p\}$ has orders $0$ and $p$ at $1$.  Only the interval bound of \cref{lem:compression} will be used.

% Adapted from an OpenAI Math manuscript. See PROVENANCE.md.
\section{An initial simplex from a flag}\label{sec:flag}

We now apply the leading-term construction to the section ring of $L$. A complete-intersection flag gives a particularly simple bound on the exponents: a simplex whose volume equals the leading coefficient of the Hilbert polynomial.  We give the complete argument for the flag construction of \cite[Section~3]{OpenAI2026Higher}.

For positive real numbers $a_1,\ldots,a_d$, write
\[
 \Delta(a_1,\ldots,a_d)
 =\left\{x\in\mathbb{R}_{\geq0}^d:
       \sum_{i=1}^d\frac{x_i}{a_i}\leq1\right\}.
\]
Its volume is $\prod_i a_i/d!$.  We use additive notation $NL$ for $L^{\otimes N}$.

\begin{proposition}\label{prop:initial-simplex}
Let $X$ be a smooth integral projective $n$-dimensional variety over an algebraically closed field, with $n\geq2$, and let $L$ be ample.  Put $V=L^n$, and choose an integer $D\geq1$ such that $DL$ is very ample. There is a graded family of finite sets $S_N\subset\mathbb{Z}_{\geq0}^n$ such that
\begin{align}
 S_N&\subset NP_0,
 &P_0&=\Delta(1/D,\ldots,1/D,D^{n-1}V),
 \label{eq:initial-simplex}\\
 \#S_N&=h^0(X,NL)
       =\frac{V}{n!}N^n+o(N^n).
 \label{eq:initial-dimension}
\end{align}
For every $N,r\geq1$ and $m\geq0$, if $M(S_N)$ separates jets through order $m$ at some $r$ distinct torus points, then $H^0(X,NL)$ separates the same jets at some $r$ distinct points of $X$.
\end{proposition}

\begin{proof}
Bertini's theorem for hyperplane sections, valid in every characteristic for a projective embedding over an algebraically closed field, gives sections $\sigma_1,\ldots,\sigma_{n-1}\in H^0(X,DL)$ whose successive zero loci form a smooth integral flag
\[
 X=X_0\supset X_1\supset\cdots\supset X_{n-1}=C_0.
\]
See \cite[Tags~\href{https://stacks.math.columbia.edu/tag/0FD6}{0FD6} and~\href{https://stacks.math.columbia.edu/tag/0G4G}{0G4G}]{Stacks} for the smoothness and integrality assertions. At each step the variety being cut has dimension at least two.  Choose $x\in C_0(k)$ and a local frame $e$ of $L$ near $x$.  The functions $z_i=\sigma_i/e^D$, $i<n$, form part of a regular system of parameters at $x$; complete them by $z_n$.  These parameters identify $\widehat{\mathcal{O}}_{X,x}$ with $k[[z_1,\ldots,z_n]]$.

Expand every section of $NL$ using the frame $e^N$, and let $S_N$ be the set of lexicographically least exponents of the nonzero expansions. The expansion map is injective: a section with zero germ is zero by integrality, and the local ring injects into its completion. \Cref{lem:least-terms} therefore gives $\#S_N=h^0(X,NL)$ and the graded property.  The Hilbert polynomial of the ample line bundle gives the last equality in \eqref{eq:initial-dimension}.  A basis with distinct least exponents, \cref{lem:exposing-weights}(i), and \cref{lem:formal-transfer} with $h=n$ prove the jet-transfer assertion.

It remains to prove the simplex bound.  Let $\alpha=(\alpha_1,\ldots,\alpha_n)\in S_N$, coming from a nonzero section $s$.  We successively divide by the equations of the flag and restrict to the next member.  The details matter because the lexicographic order first measures vanishing along an entire divisor, not just at the chosen point.

Suppose we have reached a nonzero section on $X_{i-1}$ with least exponent $(\alpha_i,\ldots,\alpha_n)$ in its formal expansion at $x$. If $\alpha_i>0$, restriction of this section to $X_i$ has zero formal expansion.  Injectivity into the completed local ring, followed by integrality of $X_i$, shows that this restriction is identically zero. The exact sequence for the effective Cartier divisor $X_i\subset X_{i-1}$ therefore allows division by $\sigma_i|_{X_{i-1}}$.  Repeating this argument $\alpha_i$ times, and then restricting to $X_i$, gives a nonzero section whose least exponent is $(\alpha_{i+1},\ldots,\alpha_n)$.  Nonvanishing follows from the nonzero coefficient of the original least term.  Each division subtracts $DL$ from the line bundle.

After these steps, the resulting nonzero section on $C_0$ belongs to
\[
 \left(N-D\sum_{i<n}\alpha_i\right)L|_{C_0}
\]
and has order exactly $\alpha_n$ at $x$.  Since $L\cdot C_0=D^{n-1}V$, comparison with the degree of its zero divisor gives
\[
 \alpha_n\leq
 \left(N-D\sum_{i<n}\alpha_i\right)D^{n-1}V.
\]
Equivalently, $D\sum_{i<n}\alpha_i+\alpha_n/(D^{n-1}V)\leq N$, which is \eqref{eq:initial-simplex}.
\end{proof}

The intercepts of $P_0$ have product $V$.  The next section changes them, two at a time, without changing this product, the numbers of exponents, or backward jet transfer.

% Adapted from an OpenAI Math manuscript. See PROVENANCE.md.
\section{Reshaping two intercepts}\label{sec:reshape}

We now change two intercepts of a simplex while preserving their product. The support sets need not fill the simplex: the construction applies to any graded family inside it.  Its algebraic step uses the two branches of a node; its remaining steps are lattice changes and coordinate compressions. This follows the planar construction of \textcite[Section~4, Proposition~4.1]{OpenAI2026Higher}, with a formal node whose branches remain distinct in characteristic two.  The nodal bound, compressions, and primitive-direction choice also appear in \textcite[Propositions~3.2 and~3.4, Lemma~4.3]{OpenAI2026Surfaces}; we give all the arguments needed here.

\begin{proposition}\label{prop:reshape}
Let $n\ge2$, let $A,B,a_3,\ldots,a_n>0$, and let $(S_N)_{N\ge1}$ be a graded family of finite subsets of $\mathbb{Z}^n$ satisfying
\[
 S_N\subset N\Delta(A,B,a_3,\ldots,a_n).
\]
Put $H=AB$ and $d=\max\{3/A,2/B\}$.  For every $w>0$ with $dw<1/16$, there is a graded family $(S'_N)_{N\ge1}$ such that
\[
 S'_N\subset N\Delta(w,H/w,a_3,\ldots,a_n),
 \qquad \#S'_N=\#S_N \quad(N\ge1).
\]
The transformation from $S_N$ to $S'_N$ has backward jet transfer in every degree.  For $n=2$, the lists $a_3,\ldots,a_n$ are empty.
\end{proposition}

We first prove the nodal bound and then choose the lattice direction that produces the prescribed intercept $w$.  Throughout this section, $H$ and $d$ have the meanings above, and we put
\begin{equation}\label{eq:reshape-constants}
 W=dH=\max\{2A,3B\},\qquad c=1/d,\qquad
 T=d^2H=\max\{9B/A,4A/B\}\ge6.
\end{equation}
In particular, $cW=H$.

\subsection{The two branches of a node}

For $t>0$, order pairs $(\alpha,\beta)\in\mathbb{Z}_{\ge0}^2$ by the lexicographic order of $(\alpha+t\beta,\beta)$, as in \cref{lem:exposing-weights}.  We denote the least exponent in this order by $\nu_t$.

\begin{lemma}\label{lem:nodal-bound}
There are formal coordinates $\xi,\eta$ at the origin of $\mathbb{A}^2_k$ with
\[
 y^2+xy-x^3=\xi\eta
\]
for which the following assertion holds.  Given $t>0$, set
\[
 a=\frac{Wt}{1+t},\qquad b=\frac{W}{1+t},\qquad
 Q=\operatorname{conv}\{(0,0),(a,0),(c,c),(0,b)\}.
\]
If $\kappa\ge0$ and $0\ne f\in k[x,y]$ has all its exponents in $\kappa\Delta(A,B)$, then its expansion in these coordinates satisfies $\nu_t(f)\in\kappa Q$.
\end{lemma}

\begin{proof}
Write $g=y^2+xy-x^3$.  The parametrization
\begin{equation}\label{eq:node-parametrization}
 x=v(v+1),\qquad y=v^2(v+1)=vx
\end{equation}
identifies the complement of $x=0$ in $g=0$ with $\operatorname{Spec}k[v,1/(v(v+1))]$: the inverse is $v=y/x$. Consequently, a polynomial vanishes under \eqref{eq:node-parametrization} if and only if it is divisible by $g$.  Indeed, vanishing implies membership in $(g)$ after inverting $x$, and $g$ is relatively prime to $x$ in the unique factorization domain $k[x,y]$.

The polynomial $Z^2+Z-x$ has simple roots $0$ and $-1$ modulo $x$. Hensel lifting gives roots $r_1,r_2\in k[[x]]$ with these respective constant terms.  Set
\[
 \eta=y-xr_1(x),\qquad \xi=y-xr_2(x).
\]
Their linear terms are $y$ and $y+x$, so they are formal coordinates, and $\xi\eta=g$.  These assertions also hold in characteristic two: the roots $0$ and $-1=1$ remain distinct and the linear terms remain independent. Along \eqref{eq:node-parametrization} near $v=0$, uniqueness of the lifted root gives $r_1(x(v))=v$, so $\eta=0$ and $\xi$ has order one in $v$. Near $v=-1$, one similarly has $\xi=0$ and $\eta$ has order one in $v+1$.

Factor $f=g^e f_0$ with $e\ge0$ and $g\nmid f_0$.  Give $x,y$ the positive weights $1/A,1/B$, and write $\deg_{A,B}$ for the maximum weight of a polynomial's monomials.  This degree is additive on products, because their highest weight forms multiply nontrivially.  Moreover,
\[
 \deg_{A,B}(g)=d,
 \qquad \deg_{A,B}(f_0)\le\kappa-ed;
\]
here $1/A+1/B\le\max\{3/A,2/B\}=d$.  Thus $ed\le\kappa$. The nonzero polynomial
\[
 F_0(v)=f_0\bigl(v(v+1),v^2(v+1)\bigr)
\]
has degree at most $W(\kappa-ed)$, since
\[
 2i+3j\le W(i/A+j/B)
\]
for every monomial $x^iy^j$ of $f_0$.  Its expansions at $v=0$ and $v=-1$ are both nonzero, since $k[v]$ injects into each completed local ring.  These expansions are exactly the restrictions of $f_0$ to the two formal branches described above.

Let $\nu_t(f_0)=(\alpha,\beta)$ and $\mu=\alpha+t\beta$. Every term of its formal expansion has weight at least $\mu$. Restricting to $\eta=0$ therefore gives order at least $\mu$ on the branch at $v=0$; restricting to $\xi=0$ gives order at least $\mu/t$ on the branch at $v=-1$.  The nonzero polynomial $F_0$ consequently satisfies
\[
 (1+1/t)\mu
 \le \ord_{v=0}(F_0)+\ord_{v=-1}(F_0)
 \le \deg(F_0)\le W(\kappa-ed).
\]
Equivalently, $(\alpha,\beta)\in(\kappa-ed)\Delta(a,b)$. Because $g=\xi\eta$, multiplicativity of least exponents gives
\[
 \nu_t(f)=(e,e)+\nu_t(f_0)
          =ed(c,c)+\nu_t(f_0)\in\kappa Q.
\]
This also covers $\kappa=0$, when $f$ is constant.
\end{proof}

\subsection{A lattice direction of prescribed height}

Choose $t>T$.  With $a,b$ as in \cref{lem:nodal-bound}, one has
\[
 a+b=W,\qquad
 \frac ca+\frac cb=\frac{(1+t)^2}{Tt}>1.
\]
Thus $Q$ is a quadrilateral with the vertices in the order displayed in that lemma, and its area is $c(a+b)/2=H/2$.  The shear $(\alpha,\beta)\mapsto(\alpha-\beta,\beta)$ sends its vertices to
\[
 (0,0),\quad(a,0),\quad(0,c),\quad(-b,b).
\]
Its vertical slice at horizontal coordinate $q$ has length
\begin{equation}\label{eq:first-slice-length}
 \begin{cases}
 c(1+q/b),&-b\le q\le0,\\
 c(1-q/a),&0\le q\le a.
 \end{cases}
\end{equation}
Compression in the second coordinate therefore places the support in the triangle
\begin{equation}\label{eq:reshape-triangle}
 P=\operatorname{conv}\{(-b,0),(a,0),(0,c)\}.
\end{equation}
Its area is again $H/2$; see \cref{fig:reshape-polygons}.

\begin{figure}[tbp]
\centering
\begin{tikzpicture}[font=\small,>=stealth]
 \begin{scope}[xscale=.9,yscale=2.4]
  \filldraw[fill=blue!9,draw=blue!60!black,thick]
   (0,0)--(4,0)--(.7,.7)--(0,.5)--cycle;
  \node[below left] at (0,0) {$0$};
  \node[below] at (4,0) {$(a,0)$};
  \node[above] at (.7,.7) {$(c,c)$};
  \node[left] at (0,.5) {$(0,b)$};
  \node at (1.65,.22) {$Q$};
 \end{scope}
 \draw[->,thick] (4.05,.8)--(5.25,.8)
  node[midway,above,align=center] {shear, then\\compress};
 \begin{scope}[shift={(6.1,0)},xscale=.9,yscale=2.4]
  \filldraw[fill=blue!9,draw=blue!60!black,thick]
   (-.5,0)--(4,0)--(0,.7)--cycle;
  \draw[densely dashed] (0,0)--(0,.7);
  \node[below] at (-.5,0) {$(-b,0)$};
  \node[below] at (4,0) {$(a,0)$};
  \node[above] at (0,.7) {$(0,c)$};
  \node at (1.6,.22) {$P$};
 \end{scope}
\end{tikzpicture}
\caption{The nodal bound $Q$ becomes the triangle $P$ after a shear and
vertical compression.  Each vertical slice of the sheared bounding polygon
is replaced by an interval of the same length starting at zero;
\eqref{eq:first-slice-length}
gives the resulting boundary.  Both polygons have area $H/2$.  The drawing
is schematic; the displayed vertex coordinates are exact.}
\label{fig:reshape-polygons}
\end{figure}

The next lemma chooses both $t$ and a primitive lattice direction.  We seek unimodular coordinates in which $P$ has vertical extent $H/w$; its area $H/2$ will then force its largest horizontal slice to have length $w$.

\begin{lemma}\label{lem:primitive-height}
If $w>0$ and $dw<1/16$, there exist $t>T$ and a primitive vector $u=(i,j)\in\mathbb{Z}^2$ such that the three vertex heights $z\mapsto\det(u,z)$ of the triangle $P$ in \eqref{eq:reshape-triangle} are distinct and their range has length $H/w$.
\end{lemma}

\begin{proof}
Put $R=H/w$ and $a_0=WT/(1+T)$.  As $t$ ranges over $(T,\infty)$, the value $a=Wt/(1+t)$ ranges over $(a_0,W)$. For $u=(i,j)$, the heights of $(-b,0),(a,0),(0,c)$ are $jb,-ja,i/d$. We seek $j>0$ and $i$ such that
\[
 ja+i/d=R,\qquad i/d>jb.
\]
The first equality turns the second inequality into $R-jW>0$. Thus we first choose $jW\le R/2$, and then use the freedom in $a$ to obtain the equality with $i$ coprime to $j$.

Take $j$ to be the largest power of two at most
\[
 \frac{R}{2W}=\frac{1}{2dw}.
\]
Then $j>1/(4dw)>4$.  The open interval
\[
 \bigl(d(R-jW),\ d(R-ja_0)\bigr)
\]
has length
\[
 dj(W-a_0)=\frac{jT}{1+T}>2,
\]
so it contains an odd integer $i$.  Since $j$ is a power of two, $(i,j)$ is primitive.  There is a unique $t>T$ such that
\[
 a=\frac{R-i/d}{j}\in(a_0,W).
\]
The resulting heights satisfy
\[
 -ja<jb<i/d,\qquad
 i/d-jb=R-jW\ge R/2>0,\qquad i/d+ja=R.
\]
This gives the required height range and distinctness.
\end{proof}

\subsection{Completion of the transformation}

\begin{proof}[Proof of \cref{prop:reshape}]
Fix $t$ and $u$ from \cref{lem:primitive-height}; these choices will be used in every degree.  Keep coordinates $3,\ldots,n$ unchanged. For each degree $N$ and each label $\gamma=(\gamma_3,\ldots,\gamma_n)$ occurring in $S_N$, put
\[
 \kappa=N-\sum_{h=3}^n\frac{\gamma_h}{a_h}\ge0.
\]
The pair exponents with this label lie in $\kappa\Delta(A,B)$. Take their polynomial span in $k[x,y]$ and replace it by its least exponents in the coordinates $\xi,\eta$ and order $\nu_t$. \Cref{lem:least-terms} preserves cardinality and gradedness: the row spaces multiply into the row with degree and label added, and the same coordinates and order are used throughout.  For fixed $N$, choose a basis with distinct least exponents in each row and apply \cref{lem:exposing-weights} to the finite union of these bases. Its common weight vector allows \cref{lem:formal-transfer} to apply to all labels at once, giving backward jet transfer.  The old points can be chosen in the dense open torus of $\mathbb{A}^2\times\mathbb{T}^{n-2}$.

\Cref{lem:nodal-bound} places the new pair exponents in $\kappa Q$. Apply the shear above and compress in the second coordinate. \Cref{lem:unimodular,lem:compression}, together with \eqref{eq:first-slice-length}, give the bound $\kappa P$. It remains to turn this triangle into $\Delta(w,H/w)$.

Complete the primitive vector $u$ to an integral basis $(u,v)$ with $\det(u,v)=1$, and use the coordinates in that basis.  This is a unimodular change, and its second coordinate is precisely $z\mapsto\det(u,z)$.  Write $q_0<q_1<q_2$ for the three vertex heights of the transformed triangle.  Its horizontal slice length is linear from zero at $q_0$ to a maximum at $q_1$, and linear back to zero at $q_2$. The maximum is $w$, since area is preserved and
\[
 \frac H2=\frac12(q_2-q_0)\,\bigl(\text{maximum slice length}\bigr),
 \qquad q_2-q_0=R=H/w.
\]
Compressing in the first coordinate therefore gives the bound
\[
 \operatorname{conv}\{(0,q_0),(w,q_1),(0,q_2)\}.
\]
Its vertical slice at first coordinate $h\in[0,w]$ has length $R(1-h/w)$.  Compression in the second coordinate gives $\Delta(w,R)$, as illustrated in \cref{fig:final-compression}.  These slice statements commute with multiplication by $\kappa$, so the resulting pair bound is $\kappa\Delta(w,R)$ for every label.  If $\kappa=0$, its only possible pair exponent is $(0,0)$, which remains unchanged at every step.

Restoring the labels gives $S'_N\subset N\Delta(w,R,a_3,\ldots,a_n)$.  Every step preserves cardinality and gradedness and has backward jet transfer, so their composition has all the asserted properties.
\end{proof}


\begin{figure}[tbp]
\centering
\begin{tikzpicture}[font=\small,>=stealth]
 \begin{scope}
  \filldraw[fill=blue!9,draw=blue!60!black,thick]
   (0,-.6)--(2,.4)--(0,1.8)--cycle;
  \draw[densely dashed] (1,-.1)--(1,1.1);
  \node[left] at (0,-.6) {$(0,q_0)$};
  \node[left] at (0,1.8) {$(0,q_2)$};
  \node[right] at (2,.4) {$(w,q_1)$};
 \end{scope}
 \draw[->,thick] (3.65,.6)--(5.15,.6)
  node[midway,above,align=center] {vertical\\compression};
 \begin{scope}[shift={(6.5,-.6)}]
  \filldraw[fill=blue!9,draw=blue!60!black,thick]
   (0,0)--(2,0)--(0,2.4)--cycle;
  \draw[densely dashed] (1,0)--(1,1.2);
  \node[below left] at (0,0) {$0$};
  \node[below] at (2,0) {$(w,0)$};
  \node[left] at (0,2.4) {$(0,R)$};
 \end{scope}
\end{tikzpicture}
\caption{After the lattice change and horizontal compression, a final
vertical compression produces the target triangle.  At first coordinate
$h$, both vertical slices have length $R(1-h/w)$, where
$R=q_2-q_0=H/w$.  Dashed segments mark corresponding slices.}
\label{fig:final-compression}
\end{figure}

% Adapted from an OpenAI Math manuscript. See PROVENANCE.md.
\section{Interior filling and interpolation}\label{sec:interpolation}

The proof needs interior lattice points of the reshaped supports, whose individual exponents are not given by the reshaping construction. We recover them by additivity across degrees, following \cite[Lemma~5.1]{OpenAI2026Higher}. This full-density argument is a special case of the semigroup approximation setting of \textcite[Theorem~1.6]{KK12}.

\begin{lemma}\label{lem:interior-filling}
Let $P\subset\mathbb{R}^n$, $n\ge1$, be a compact convex set with nonempty interior. Suppose finite sets $F_N\subset NP\cap\mathbb{Z}^n$, indexed by integers $N\ge1$, satisfy
\[
 F_N+F_{N'}\subset F_{N+N'}
 \quad\text{and}\quad
 \#F_N=\vol(P)N^n+o(N^n).
\]
Then for every compact set $K\subset\operatorname{int}(P)$ there is an integer $N_K$ such that
\[
 NK\cap\mathbb{Z}^n\subset F_N\qquad(N\ge N_K).
\]
\end{lemma}

\begin{proof}
A convex body has boundary of measure zero. The lattice-point asymptotic therefore gives $\#((NP\cap\mathbb Z^n)\setminus F_N)=o(N^n)$. Split $N=N_1+N_2$ with $N_1=\lfloor N/2\rfloor$. Choose $\rho>0$ such that the closed $\rho$-neighborhood of $K$ lies in $P$. For $y\in NK\cap\mathbb Z^n$, the intersection $N_1P\cap(y-N_2P)$ contains the ball of radius $\rho N_1$ about $(N_1/N)y$. By inscribing a cube, the ball contains $\gg N^n$ lattice points uniformly in $y$. Of these candidates for $y_1$, only $o(N^n)$ can have $y_1\notin F_{N_1}$ or $y-y_1\notin F_{N_2}$. A remaining decomposition gives $y\in F_{N_1}+F_{N_2}\subset F_N$.
\end{proof}

In the stretched simplex, the admissible powers of the last coordinate interpolate the coefficient of each transverse monomial independently, as in \cite[Lemma~5.2]{OpenAI2026Higher}.

\begin{lemma}\label{lem:elongated-jets}
Let $n,r\ge1$ and $m\ge0$ be integers, and let $k$ be a field containing distinct nonzero elements $c_1,\ldots,c_r$.  Set
\begin{equation}\label{eq:interpolation-support}
 E_m=\left\{(\beta,b)\in\mathbb{Z}_{\ge0}^{n-1}\times\mathbb{Z}_{\ge0}:
          |\beta|+\frac br<m+1\right\},
 \qquad |\beta|=\sum_{j=1}^{n-1}\beta_j.
\end{equation}
Then evaluation on jets through order $m$ at the $r$ points
\[
 q_i=(1,\ldots,1,c_i)\in\mathbb{T}^n(k)
\]
is an isomorphism from $M(E_m)$ to the direct sum of the jet spaces. For $n=1$, the tuple $\beta$ is empty and $|\beta|=0$.
\end{lemma}

\begin{proof}
Write the torus coordinates as $(x',Y)=(x_1,\ldots,x_{n-1},Y)$ and put $X=x'-(1,\ldots,1)$. Translation from $x'$ to $X$ preserves the space of polynomials with the indicated support. In fact, the coefficient space of $Y^b$ consists of the polynomials of total degree less than $m+1-b/r$ in the first $n-1$ variables. Consequently every element of the space has a unique expression
\[
 \sum_{|\alpha|\le m}X^\alpha f_\alpha(Y),
 \qquad
 \deg f_\alpha<r(m+1-|\alpha|).
\]
At $q_i$ the local parameters are $X$ and $Y-c_i$. For each $\alpha$ the jet records the remainder of $f_\alpha$ modulo $(Y-c_i)^{m+1-|\alpha|}$. As $i$ varies, these powers are pairwise comaximal. The Chinese remainder theorem gives a bijection between the simultaneous remainders and the polynomials of degree less than the degree $r(m+1-|\alpha|)$ of their product. Taking this bijection for each $\alpha$ proves the assertion.
\end{proof}

For large $r$, successive reshapings put a translated copy of this interpolation support in the interior.

\begin{proposition}\label{prop:generic-jets}
Let $k$ be an algebraically closed field of positive characteristic, let $X/k$ be a smooth integral projective variety of dimension $n\ge2$, and let $L$ be ample.  Put $V=L^n$.  There is an integer $r_0=r_0(X,L)$ such that, for every integer $r\ge r_0$ and every real number
\[
 0<\theta<\left(\frac Vr\right)^{1/n},
\]
there is an integer $N_0=N_0(r,\theta)$ such that, for every $N\ge N_0$, the sections of $NL$ separate jets through order $\lfloor N\theta\rfloor$ at some ordered $r$-tuple of distinct $k$-points, which may depend on $N$. For each such $N$, the same jet separation holds after base change to $K_r$ at the geometric generic tuple $p_1,\ldots,p_r$.
\end{proposition}

\begin{proof}
Take the graded family $S_N$ and simplex $P_0=\Delta(a_1,\ldots,a_n)$ from \cref{prop:initial-simplex}. Thus $\prod_j a_j=V$, the support $S_N$ lies in $NP_0$, and
\[
 \#S_N=\frac{V}{n!}N^n+o(N^n).
\]
Let $\ell=\min_j a_j$, and choose $r_0$ so large that
\[
 w:=\left(\frac Vr\right)^{1/n}<\frac{\ell}{100}
 \qquad(r\ge r_0).
\]
Fix one such $r$.  Reshape the first two intercepts to $(w,a_1a_2/w)$, then reshape the second and third intercepts, and continue through the $(n-1)$st and $n$th intercepts.  At every step, both intercepts $A,B$ being reshaped are at least $\ell$: untouched intercepts have this property, and the carried intercept $AB/w$ is larger than $A$ because $B>w$. Consequently
\[
 w\max(3/A,2/B)\le\frac{3w}{\ell}<\frac1{16},
\]
as required by \cref{prop:reshape}.  The product of all intercepts remains $V$.  The resulting graded family $F_N$ therefore satisfies
\begin{equation}\label{eq:elongated-full-density}
 \begin{aligned}
 &F_N\subset NP_*\cap\mathbb{Z}^n,
 \qquad P_*:=\Delta(w,\ldots,w,rw),\\
 &\#F_N=\frac{V}{n!}N^n+o(N^n)
        =\vol(P_*)N^n+o(N^n).
 \end{aligned}
\end{equation}
Here the last intercept is $V/w^{n-1}=rw$. The transformations preserve backward jet transfer, so jet separation by $M(F_N)$ will imply jet separation by $H^0(X,NL)$.

Fix $0<\theta<w$.  To apply \cref{lem:interior-filling}, we contract $P_*$ and then translate it away from the coordinate faces. Choose $\lambda$ with $\theta/w<\lambda<1$, and choose a vector $v\in\mathbb{R}_{>0}^n$ small enough that
\[
 v+\lambda P_*\subset\operatorname{int}(P_*).
\]
For example, it suffices that $\sum_{j<n}v_j/w+v_n/(rw)<1-\lambda$. Choose vectors $v_N\in N^{-1}\mathbb{Z}^n$ tending to $v$. For all sufficiently large $N$, the sets $v_N+\lambda P_*$ lie in a common compact subset $K$ of $\operatorname{int}(P_*)$.  Applying \cref{lem:interior-filling} to \eqref{eq:elongated-full-density} gives
\begin{equation}\label{eq:translated-inner-support}
 Nv_N+\bigl(N\lambda P_*\cap\mathbb{Z}^n\bigr)\subset F_N
\end{equation}
for all sufficiently large $N$.

Put $m=\lfloor N\theta\rfloor$.  Since $\theta<\lambda w$, we have $m+1\le N\lambda w$ for all sufficiently large $N$. The defining inequality \eqref{eq:interpolation-support} therefore gives
\[
 E_m\subset N\lambda P_*\cap\mathbb{Z}^n.
\]
Choose distinct $c_1,\ldots,c_r\in k^\times$ and apply \cref{lem:elongated-jets}.  Multiplication by the Laurent monomial with exponent $Nv_N\in\mathbb{Z}^n$ is multiplication by a unit on the torus, hence induces an automorphism on every jet space.  Thus \eqref{eq:translated-inner-support} implies that $M(F_N)$ separates jets through order $m$ at these $r$ points.  Backward jet transfer through all reshapings and \cref{prop:initial-simplex} gives the asserted jet separation at an ordered tuple on $X$.

For each fixed $r$, $\theta$, and sufficiently large $N$, the locus in $U_r$ where this evaluation map is surjective is a nonempty open set, by the openness of jet separation established in \cref{sec:transfer}. The variety $U_r$ is integral, so its generic point belongs to that open set. Global sections of $NL$ commute with field extension, as do principal parts. Base change of the evaluation map at the generic point to $K_r$ therefore gives the jet evaluation of $H^0(X_{K_r},NL_{K_r})$ at $p_1,\ldots,p_r$.  The same geometric generic tuple therefore satisfies every required jet-separation statement.
\end{proof}

% Adapted from an OpenAI Math manuscript. See PROVENANCE.md.
\section{From jets to the Seshadri constant}\label{sec:positivity}

The interpolation statement now gives the lower bound for the Seshadri constant.  We first spell out the passage from jets to intersection numbers, including at singular points of a curve, and then recall the matching volume bound.  The relation between asymptotic jet separation and Seshadri constants is classical; see \cite[Theorem~6.4]{Demailly1992} in the complex setting and \cite[Section~2, Equation~(2.3)]{MustataSchwede2014} in arbitrary characteristic. The arguments below work in arbitrary characteristic.

\begin{lemma}\label{lem:jet-curve-bound}
Let $Y$ be a smooth integral projective variety over an algebraically closed field, let $A$ be a line bundle on $Y$, and let $q_1,\ldots,q_r$ be distinct points.  Suppose that $H^0(Y,A)$ separates jets through order $m\geq 1$ at these points.  Then every integral curve $C\subset Y$ meeting at least one of them satisfies
\[
 A\cdot C\geq m\sum_{i=1}^r\mult_{q_i}C,
\]
where the multiplicity is zero at a point outside $C$.
\end{lemma}

\begin{proof}
Choose such a point $q_j$.  Write $R=\mathcal{O}_{C,q_j}$ and let $\mathfrak{n}$ be its maximal ideal.  The vector space $\mathfrak{n}^m/\mathfrak{n}^{m+1}$ is nonzero: otherwise Nakayama's lemma would give $\mathfrak{n}^m=0$, contrary to $R$ being a one-dimensional local domain. In a local frame of $A$, lift a nonzero class in this vector space to an ambient jet of order $m$ at $q_j$.  This lift can be chosen in the image of $\mathfrak{m}_{Y,q_j}^m$.  Prescribe that jet at $q_j$ and the zero jet at all other marked points.  Jet separation produces a section $\sigma\in H^0(Y,A)$ which vanishes to order at least $m$ at every $q_i$ and whose restriction to $C$ is nonzero.

For completeness, the local multiplicity estimate used here is
\begin{equation}\label{eq:local-zero-length}
 \operatorname{length}_R(R/(f))\geq m\,e(R)
 \qquad(0\ne f\in\mathfrak{n}^m),
\end{equation}
where $e(R)$ denotes the Hilbert--Samuel multiplicity of a one-dimensional Noetherian local domain $(R,\mathfrak{n})$; see \cite[Definition~43.15.1, Tag~\href{https://stacks.math.columbia.edu/tag/0AZV}{0AZV}]{Stacks}.  Indeed, for positive integers $t$, the nonzero divisor $f$ gives
\[
 t\operatorname{length}_R(R/(f))
 =\operatorname{length}_R(R/(f^t))
 \geq\operatorname{length}_R(R/\mathfrak{n}^{mt})
 =mt\,e(R)+O(1).
\]
Divide by $t$ and let $t$ tend to infinity.  Applying \eqref{eq:local-zero-length} to the local equations of $\sigma|_C$ at the marked points proves the inequality: the sum of all local zero lengths of a nonzero section of $A|_C$ is $\deg(A|_C)=A\cdot C$.
\end{proof}

\begin{lemma}\label{lem:volume-bound}
Let $Y$ be a smooth integral projective variety of dimension $d\geq2$ over an algebraically closed field, let $A$ be ample, and let $q_1,\ldots,q_r$ be distinct points.  Then
\[
 \inf_{C\cap\{q_1,\ldots,q_r\}\ne\varnothing}
 \frac{A\cdot C}{\sum_i\mult_{q_i}C}
 \leq \left(\frac{A^d}{r}\right)^{1/d},
\]
where the infimum is over integral curves in $Y$.
\end{lemma}

\begin{proof}
Denote the infimum by $\varepsilon$.  If $\varepsilon=0$, the conclusion is immediate.  Otherwise, blow up the marked points:
\[
 \pi:\widetilde Y\longrightarrow Y,
 \qquad E_1,\ldots,E_r\subset\widetilde Y.
\]
For each rational number $0\leq s<\varepsilon$, the divisor class
\[
 D_s=\pi^*A-s\sum_i E_i
\]
is nef.  To check this on curves, first observe that $D_s|_{E_i}=s\mathcal{O}_{\mathbb{P}^{d-1}}(1)$, which has nonnegative degree on every curve in $E_i$.  Every other integral curve in $\widetilde Y$ is the strict transform $\widetilde C$ of an integral curve $C\subset Y$, and
\begin{equation}\label{eq:exceptional-multiplicity}
 E_i\cdot\widetilde C=\mult_{q_i}C.
\end{equation}
Thus $D_s\cdot\widetilde C\geq0$ by the definition of $\varepsilon$; for curves missing all marked points this follows from the ampleness of $A$.

Here is a local justification of \eqref{eq:exceptional-multiplicity}, which also covers singular curves.  If $q_i\in C$, the Rees algebra construction identifies $\widetilde C$ with the blow-up of $C$ at the marked points. Indeed, the ambient Rees algebra surjects onto the Rees algebra of the restricted ideal.  The latter is integral, and its Proj agrees with $C$ away from the centers, so its closed image is exactly the strict transform. For $R=\mathcal{O}_{C,q_i}$ with maximal ideal $\mathfrak{n}$, its exceptional fiber is
\[
 \operatorname{Proj}\bigl(\operatorname{gr}_{\mathfrak{n}}R\bigr).
\]
It is the effective Cartier divisor cut out on $\widetilde C$ by $E_i$. Its length is $e(R)$: the Hilbert function $\dim_{k(q_i)}(\mathfrak{n}^t/\mathfrak{n}^{t+1})$ equals $e(R)$ for all sufficiently large $t$, and its eventual value is the length of this zero-dimensional projective scheme.  This proves \eqref{eq:exceptional-multiplicity}; if $q_i\notin C$, both sides vanish.

A nef class on a projective variety is a limit of ample classes and therefore has nonnegative top self-intersection.  The exceptional divisors are disjoint, $\pi^*A$ restricts trivially to each of them, and $E_i^d=(-1)^{d-1}$.  Consequently
\[
 0\leq D_s^d=A^d-rs^d.
\]
If $\varepsilon>(A^d/r)^{1/d}$, even if $\varepsilon=+\infty$, there is a rational number $s$ strictly between $(A^d/r)^{1/d}$ and $\varepsilon$. The displayed inequality excludes such an $s$, proving the bound.
\end{proof}

\begin{proof}[Proof of \cref{thm:main}]
Fix $r\ge r_0$ as in \cref{prop:generic-jets} and put $w=(L^n/r)^{1/n}$. For $0<\theta<w$, that proposition separates jets of order $\lfloor N\theta\rfloor$ at the geometric generic tuple over $K_r$, for all large $N$. By \cref{lem:jet-curve-bound}, $\varepsilon_r(X,L)\ge\lfloor N\theta\rfloor/N$. Letting $N\to\infty$ and then $\theta\uparrow w$ gives the lower bound. The upper bound is \cref{lem:volume-bound}, since $(L_{K_r})^n=L^n$. On the blow-up of the geometric generic tuple, the same curve calculation shows that $\pi^*L_{K_r}-w\sum_i E_i$ is nef, by taking rational $s\uparrow w$. Its top self-intersection is $L^n-rw^n=0$.
\end{proof}


\printbibliography

\end{document}
